Rotation smoothing is a fundamental problem in robotics and computer vision, where sequences of noisy rotation measurements must be denoised while preserving smooth motion. Applications span camera pose estimation, attitude control, inertial navigation, and 3D reconstruction, where rotation data comes from sensors such as IMUs, visual odometry, or structure-from-motion pipelines.

Traditional approaches to rotation smoothing typically assume Gaussian noise models and minimize energy functions penalizing angular velocity and acceleration. However, many practical scenarios involve bounded-error constraints where sensor specifications, calibration uncertainty, or safety requirements mandate that smoothed estimates remain within known geodesic bounds from the measurements. This set-membership formulation, first introduced in optimal estimation theory for linear systems, has received limited attention for nonlinear manifolds such as SO(3).

The challenges in SO(3) rotation smoothing are threefold. First, the nonlinear manifold structure precludes direct application of linear smoothing techniques. Second, hard tube constraints make the optimization problem non-convex and computationally challenging. Third, the requirement to maintain feasibility with respect to all constraints demands careful handling of the manifold geometry during optimization.

In this paper, we present a bounded-noise tube smoothing algorithm for SO(3) that addresses these challenges through sequential convexification combined with a structured Alternating Direction Method of Multipliers (ADMM) solver. Our approach enforces that each smoothed rotation lies within a geodesic ball (tube) of radius $\epsilon_i$ around the corresponding measurement, where $\epsilon_i$ represents a known upper bound on the estimation error rather than a tuned parameter.

The key innovation is reformulating the problem as a sequence of convex SOCP subproblems, each solvable via a custom ADMM implementation that exploits the problem's block-banded Hessian structure. This yields three main advantages: (1) global feasibility is maintained through tube constraints; (2) computational efficiency scales near-linearly with problem size; (3) theoretical guarantees include convergence rate bounds and optimality condition verification.

We make three primary contributions:

\textbf{1. Set-membership tube smoothing on SO(3).} We formulate rotation smoothing as a bounded-noise estimation problem where each smoothed output must lie within a measurement-derived tube. This contrasts with traditional Gaussian assumptions and provides a rigorous framework for safety-critical applications where error bounds are specified by sensor specifications or calibration data.

\textbf{2. Sequential convexification with trust-region management.} We employ local retractions on the manifold to update rotations, ensuring that each step remains within the validity region of the linearized tube constraint. The trust-region mechanism dynamically adjusts based on actual-to-predicted improvement ratios, providing robustness across varying problem conditions.

\textbf{3. Structured ADMM solver with theoretical enhancements.} We design a custom ADMM solver that exploits the block-banded structure of the smoothing Hessian, enabling efficient sparse linear solves that scale to large problem sizes ($M \geq 10^5$). The implementation includes warm-starting strategies, adaptive trust-region parameters, Karush-Kuhn-Tucker optimality condition verification, and convergence rate analysis.

We demonstrate the effectiveness of our approach on both synthetic bounded-noise problems and real sensor data. Experiments show that our method maintains tube constraint feasibility while achieving smooth trajectories comparable to unconstrained baselines. The structured ADMM solver provides 15--30\% speedup over general SOCP solvers, with near-linear scaling behavior demonstrated up to problem sizes of $M = 10^5$ rotation samples.

The remainder of this paper is organized as follows. Section~\ref{sec:related} reviews related work in rotation smoothing and set-membership estimation. Section~\ref{sec:methods} presents our problem formulation, sequential convexification framework, and structured ADMM solver. Section~\ref{sec:theory} provides theoretical analysis including convergence guarantees and optimality conditions. Section~\ref{sec:experiments} describes our experimental setup and benchmarks. Section~\ref{sec:results} presents evaluation results and discussion. Section~\ref{sec:conclusion} concludes with limitations and future directions.
